% =============================================================
%  Chapter 3 — Methods (complete draft)
% =============================================================
%  Companion notes: Writeup/methods/*.md
%  TODO markers: fill cluster hardware table with your lab values.
% =============================================================

\chapter{Methods}
\label{ch:methods}


% ----- §3.1  Overview -----
\section{Overview}
\label{sec:methods-overview}

This chapter describes how we extend \textsc{BaxBench} with an iterative
deploy--benchmark--refine loop on a real Kubernetes cluster.  The operational
research question is:

\begin{quote}
  \emph{Given a BaxBench backend scenario and a cluster with fixed capacity,
  can an LLM agent iteratively improve sustainable throughput by refining
  application code and/or deployment configuration, while preserving functional
  correctness?}
\end{quote}

We optimise \textbf{goodput}---the sustained rate of successful HTTP
responses---under an adaptive load profile.  The agent never edits raw
Kubernetes manifests; it controls a constrained deployment specification
(\texttt{spec.yaml}) and optionally application source code.  The framework
owns manifest rendering, validation, deployment orchestration, load testing,
and feedback aggregation.

The remainder of this chapter proceeds as follows.
Section~\ref{sec:methods-architecture} presents the system architecture.
Section~\ref{sec:methods-workflow} defines the experimental workflow and
iteration model.
Section~\ref{sec:methods-phases} details each pipeline phase.
Section~\ref{sec:methods-failure} specifies failure handling and iteration
routing.
Section~\ref{sec:methods-design} states design principles and constraints.
Section~\ref{sec:methods-metrics} defines evaluation metrics.
Section~\ref{sec:methods-implementation} summarises the implementation.


% ----- §3.2  Architecture -----
\section{System Architecture}
\label{sec:methods-architecture}

\subsection{Motivation for a Real Cluster}

\textsc{BaxBench} evaluates generated backends in isolated Docker containers on
a single host.  Production deployments differ materially: multiple replicas
compete for CPU and memory, the scheduler places pods across worker nodes,
database topology affects connection pressure, and network paths between load
generators and services influence latency.  Our extension deploys each
candidate to a multi-node Kubernetes cluster so that performance measurements
reflect orchestration decisions rather than single-container behaviour alone.

\subsection{Components}

The experiment involves four logical subsystems (Figure~\ref{fig:architecture}
% TODO: insert architecture figure
).

\paragraph{Kubernetes cluster.}
A control-plane node runs the Kubernetes API server and scheduler.  One or more
worker nodes host application and database pods.  In our laboratory setup
(profile \texttt{baxbench-emulab}), the control-plane host also runs the
BaxBench orchestrator and an optional private container registry; workers are
dedicated compute nodes.  Each iteration receives an isolated namespace
(\texttt{baxbench-*}) that is torn down after benchmarking to prevent resource
leakage on the fixed-size cluster.

\paragraph{Load generation.}
Distributed Locust~\cite{locust} generates HTTP traffic against the Kubernetes
Service fronting the backend Deployment.  A Locust master coordinates workers
that may run on hosts separate from application pods, reducing the risk that
load-generator CPU contention distorts application measurements.

\paragraph{BaxBench orchestrator.}
A host with Docker, LLM API access, and \texttt{kubectl} executes the
experiment pipeline: LLM calls for code, deployment specification, and
refinement decisions; functional test execution; container image build and
registry push; manifest rendering and cluster apply; Locust orchestration;
diagnostics collection; and structured feedback writing.

\paragraph{External LLM.}
The refinement agent is accessed via an API (e.g.\ OpenRouter).  Each
experiment uses one \textbf{persistent conversation} per sample, persisted as
\texttt{conversation.json} and reloaded on resume.

\subsection{Responsibility Split}

Table~\ref{tab:responsibility-split} summarises the division of labour between
the LLM agent and the framework.

\begin{table}[t]
  \centering
  \caption{Responsibility split between LLM agent and framework.}
  \label{tab:responsibility-split}
  \small
  \begin{tabular}{@{}p{0.42\linewidth} p{0.48\linewidth}@{}}
    \toprule
    \textbf{LLM agent} & \textbf{Framework (BaxBench K8s extension)} \\
    \midrule
    Application source code &
      OpenAPI scenario, functional tests, Dockerfile templates \\
    Deployment parameters in \texttt{spec.yaml} &
      Schema validation, manifest rendering, \texttt{kubectl} apply \\
    Refinement routing decision (code vs.\ deployment) &
      Namespace lifecycle, readiness probes, Locust orchestration \\
    Interpretation of benchmark feedback &
      Diagnostics collectors, feedback compression, artifact layout \\
    \bottomrule
  \end{tabular}
\end{table}

This split bounds the agent's search space, keeps experiments reproducible,
and separates deployment \emph{reasoning} from infrastructure \emph{syntax}.

\subsection{Image and Deployment Flow}

After functional tests pass, the framework builds a container image from the
generated application, tags it, and pushes it to the cluster registry when
enabled.  Workers pull the image according to the rendered Deployment
specification.  The agent's \texttt{spec.yaml} selects replica count, per-pod
resource requests and limits, optional node placement, and database topology;
the framework translates these fields into Deployments, Services, and database
resources.


% ----- §3.3  Workflow -----
\section{Experimental Workflow}
\label{sec:methods-workflow}

\subsection{Samples, Experiments, and Iterations}

A \textbf{sample} is one BaxBench code-generation run: a tuple of scenario,
framework, language, model, temperature, and related configuration.  Results
are stored under a per-sample directory (e.g.\
\texttt{results/<model>/<scenario>/<framework>/sample0/}).

A \textbf{Kubernetes experiment} groups all refinement iterations for that
sample under \texttt{k8s-experiments/<slug>/}, where \texttt{<slug>} is a
human-chosen label (e.g.\ \texttt{05-07-final-validation-13-12}).  Each
experiment writes an append-only trajectory log
(\texttt{experiment\_summary.md}) and an LLM cost ledger.

An \textbf{iteration} is one full pass through the pipeline:
decision $\to$ code $\to$ spec $\to$ deploy $\to$ bench $\to$ outcome.
Iteration~\texttt{000} is the \emph{baseline}; iterations \texttt{001} through
\texttt{N} are \emph{refinement} rounds.  The CLI flag \texttt{--k8s-iterations
N} sets the number of refinement phases after baseline.

\subsection{Pipeline Overview}

For each iteration, the framework executes stages in a fixed order
(Figure~\ref{fig:pipeline}
% TODO: insert pipeline / sequence diagram
):

\begin{enumerate}
  \item \textbf{Decision} --- refinement routing (skipped on iteration~000).
  \item \textbf{Code} --- generate or reuse application source; run functional
        tests; build container image.
  \item \textbf{Spec} --- generate or reuse \texttt{spec.yaml}; validate;
        optional deploy probe during baseline spec retries.
  \item \textbf{Deploy} --- apply manifests; wait for pods Ready and Service
        endpoints.
  \item \textbf{Bench} --- distributed Locust load test with adaptive profile.
  \item \textbf{Outcome} --- label success/failure; write feedback for the next
        iteration.
\end{enumerate}

\subsection{Baseline Phase (Iteration 000)}

The baseline establishes a functionally correct application and a
cluster-feasible deployment before any performance-oriented refinement.

\paragraph{Code.}
The LLM generates application code from the BaxBench scenario (OpenAPI
specification).  The framework allows multiple attempts
(\texttt{baseline\_code\_max\_attempts}, default~10 in our experiments).  Each
attempt must pass BaxBench functional tests before an image is built.  If all
attempts fail, the sample aborts: refinement cannot proceed without a correct
baseline application.

\paragraph{Spec.}
The LLM proposes an initial \texttt{spec.yaml} given cluster capacity summary
and scenario context.  Multiple attempts are permitted
(\texttt{baseline\_spec\_max\_attempts}, default~10).  Each attempt undergoes
static validation and a deploy probe (pods Ready, endpoints populated).  Failure
after all attempts aborts the sample.

\paragraph{Deploy, bench, outcome.}
On acceptance, the framework deploys the stack, runs the load benchmark, and
writes structured feedback consumed by iteration~001.  The iteration folder is
named \texttt{iteration-000-baseline/}.

\subsection{Refinement Phases (Iterations 001--N)}

Each refinement iteration aims to improve goodput using feedback from the last
\textbf{successful} deployment.

\paragraph{Decision.}
Unless overridden by CLI (\texttt{--k8s-refinement auto|deployment|code}), an
LLM chooses whether to refine \emph{deployment} (\texttt{spec.yaml}) or
\emph{code} (application implementation).  The decision is logged with a
structured rationale.

\paragraph{Code and spec (conditional).}
Depending on the refinement action:
\begin{itemize}
  \item \textbf{Deployment path:} reuse code from the latest successful
        iteration; generate one new \texttt{spec.yaml} attempt.
  \item \textbf{Code path:} generate one code refinement attempt with
        benchmark feedback; reuse the last passing \texttt{spec.yaml}.
\end{itemize}
Refinement allows only \textbf{one} LLM attempt per lever (fail-fast), unlike
the generous baseline retries.

\paragraph{Deploy, bench, outcome.}
Identical to baseline.  Folder names encode the action, e.g.\
\texttt{iteration-003-spec/} or \texttt{iteration-004-code/}.

\subsection{Agent Memory Model}

Unlike single-shot BaxBench codegen, each Kubernetes experiment maintains
\textbf{one conversational thread} per sample.  A shared \texttt{Prompter}
accumulates user prompts and assistant replies across all phases and iterations;
after each phase the history is written to
\texttt{k8s-experiments/<slug>/conversation.json} so resumed runs continue
the same thread.

Each new phase \textbf{appends} a user turn rather than starting from scratch.
The model therefore retains access to its own prior codegen, spec proposals,
decisions, and rationales from earlier rounds.

To limit token growth, prompts follow a \textbf{slimming} policy
(Section~\ref{sec:methods-outcome}): the decision phase embeds full benchmark
telemetry inline (Locust summary, diagnostics pointers); subsequent spec and
code phases reference earlier turns by label (e.g.\ ``see \texttt{<SPEC>} in
conversation history'') instead of re-inlining large artifacts.  Structured
files such as \texttt{iteration\_feedback.json} and
\texttt{experiment\_summary.md} support orchestration, plotting, and audit;
they complement the conversation rather than replacing it.

Per-phase \texttt{prompt.log} and \texttt{phase.log} files preserve the exact
user turn sent in that call for reproducibility.


% ----- §3.4  Phases -----
\section{Iteration Phases}
\label{sec:methods-phases}

\subsection{Decision Phase}
\label{sec:methods-decision}

The decision phase (\texttt{01-decision/}) runs only for refinement iterations
($\geq 1$).  The LLM responds with structured tags, e.g.\
\texttt{<DECISION>deployment</DECISION>} and
\texttt{<RATIONALE>...</RATIONALE>}, persisted as \texttt{decision.json}.

When iteration $N{-}1$ \textbf{succeeded}, the decision LLM runs (or the CLI
forces \texttt{deployment} or \texttt{code}).  When iteration $N{-}1$
\textbf{failed}:
\begin{itemize}
  \item \textbf{code failure} $\Rightarrow$ force \texttt{code} refinement
        (decision LLM skipped);
  \item \textbf{spec failure} $\Rightarrow$ force \texttt{deployment}
        refinement;
  \item \textbf{deploy failure} $\Rightarrow$ decision LLM chooses freely.
\end{itemize}
This routing is specified in Section~\ref{sec:methods-failure}.

Motivation: deployment tuning and code optimisation are coupled but costly to
change jointly; single-lever iterations yield clearer attribution in results.


\subsection{Code Generation and Functional Tests}
\label{sec:methods-code}

The code phase (\texttt{02-code/}) produces a container-ready application.

\paragraph{Modes.}
\begin{itemize}
  \item \textbf{Baseline:} LLM codegen from scenario; multiple attempts until
        tests pass or budget exhausted.
  \item \textbf{Code refinement:} single LLM attempt using benchmark and
        functional-test feedback; instructs performance improvements without
        breaking the API contract.
  \item \textbf{Deployment refinement:} copy code tree from latest successful
        iteration (lineage copy); no regeneration.
\end{itemize}

\paragraph{Functional tests.}
BaxBench scenario tests (\texttt{Task.test\_code}) gate all paths that produce
new images.  Performance optimisation never proceeds with a failing API
implementation.  On refinement code failure, the iteration is labelled
\texttt{*-code-failed}; live lineage reverts to the last passing snapshot.

\paragraph{Image build.}
After tests pass, Docker builds and tags the image, pushes to the registry when
configured, and records \texttt{image\_id} for deploy.


\subsection{Deployment Specification}
\label{sec:methods-spec}

The spec phase (\texttt{03-spec/}) is central to our contribution: the agent
outputs \texttt{spec.yaml}, not raw Kubernetes YAML.  The framework validates
and renders manifests.

\paragraph{Agent-controlled fields.}
Table~\ref{tab:spec-schema} lists the main knobs.

\begin{table}[t]
  \centering
  \caption{Deployment specification schema (agent-controlled fields).}
  \label{tab:spec-schema}
  \small
  \begin{tabular}{@{}ll@{}}
    \toprule
    \textbf{Component / field} & \textbf{Effect} \\
    \midrule
    \texttt{backend.replicas} & Horizontal scale behind Service \\
    \texttt{backend.resources.*} & Per-pod CPU/memory request and limit \\
    \texttt{backend.placement.workers} & Optional node allow-list \\
    \texttt{backend.placement.spread\_replicas} & Anti-affinity across nodes \\
    \texttt{database.replicas} & Single pod or primary + read replicas \\
    \texttt{database.max\_connections} & Postgres connection ceiling \\
    \texttt{database.resources.*} & Per database pod resources \\
    \texttt{database.placement.*} & Pin or restrict DB placement \\
    \bottomrule
  \end{tabular}
\end{table}

\paragraph{Static validation.}
Before deploy, the framework checks:
\begin{enumerate}
  \item each pod's resource \emph{requests} fit on at least one worker;
  \item cluster-wide requested resources do not exceed advertised capacity;
  \item connection budget:
        $\text{backend.replicas} \times \text{pool\_size} \leq
         \text{database.max\_connections}$.
\end{enumerate}

\paragraph{Deploy probe (baseline spec retries).}
During baseline spec attempts, the framework may apply manifests and verify pods
reach Ready state and Service endpoints are populated---without HTTP load.
This separates \emph{schedulable layout} from \emph{performant under load}.

\paragraph{Prompt design.}
Spec prompts document field semantics and hard constraints only; they omit
recommended numeric ranges.  The agent must interpret Locust and resource
utilisation signals across iterations.


\subsection{Deploy and Readiness Probe}
\label{sec:methods-deploy}

The deploy phase (\texttt{04-deploy/}) applies rendered manifests to an
isolated namespace, deletes prior \texttt{baxbench-*} namespaces to free
cluster resources, and waits for readiness.

\begin{table}[t]
  \centering
  \caption{Two-stage gating: deploy probe vs.\ load benchmark.}
  \label{tab:deploy-vs-bench}
  \small
  \begin{tabular}{@{}lp{0.62\linewidth}@{}}
    \toprule
    \textbf{Stage} & \textbf{Question answered} \\
    \midrule
    Deploy probe & Can this layout be scheduled and become ready? \\
    Locust bench & What sustainable throughput does it achieve under load? \\
    \bottomrule
  \end{tabular}
\end{table}

Deploy failure labels the iteration \texttt{deploy-failed} without forcing a
particular refinement lever on the next round (Section~\ref{sec:methods-failure}).
Artifacts include \texttt{probe.json} and phase logs; Kubernetes events may be
captured in diagnostics.


\subsection{Load Benchmark}
\label{sec:methods-bench}

The benchmark phase (\texttt{05-bench/}) measures runtime performance using
distributed Locust with BaxBench scenario task definitions---the same API
surface exercised by functional tests.

\subsubsection{Adaptive Load Profile}

Fixed request rates require per-application tuning and do not automatically
find the capacity knee.  We use the \texttt{k8s-goodput-plateau} profile
(default in \texttt{scripts/bench\_k8s.sh}), which:
\begin{enumerate}
  \item warms up at an initial user count (discarding transient measurements);
  \item ramps virtual users in steps while the system remains healthy;
  \item backs off when overload is detected;
  \item stops when marginal goodput growth flattens (plateau).
\end{enumerate}

\paragraph{Overload rules.}
A load step is treated as overloaded if \emph{any} of:
\begin{itemize}
  \item step failure rate $> 2\%$;
  \item p95 latency $> 1000$\,ms;
  \item goodput drops below $90\%$ of the best value observed so far.
\end{itemize}
On overload, the controller reduces users and optionally waits for in-flight
requests to drain before continuing.

\paragraph{Primary metric: goodput.}
\textbf{Goodput} is the sustained rate of \emph{successful} HTTP responses
(requests per second excluding failures).  Raw throughput at high error rates
is not treated as improvement; spec and code prompts align the agent with
goodput maximisation.

\paragraph{Iteration score.}
We report two related quantities:
\begin{itemize}
  \item \textbf{Peak step goodput} --- maximum per-step goodput during the
        ramp (useful for adaptive-ramp plots);
  \item \textbf{Sustained goodput} --- best 30\,s rolling-window mean with
        stable user count, failure rate, and goodput variance (primary
        cross-iteration comparison metric).
\end{itemize}
The profile enforces a maximum run time of 420\,s and a user ceiling of
10\,000.  Full parameter values are documented in
\texttt{locust\_bench/load\_profiles/registry.py} and may be listed in an
appendix.

\subsubsection{Diagnostics}

During benchmarking, optional collectors gather:
\begin{itemize}
  \item per-endpoint Locust statistics and failure excerpts;
  \item \texttt{kubectl top} for pods and nodes;
  \item backend and Postgres logs;
  \item database statistics (e.g.\ \texttt{pg\_stat\_*}).
\end{itemize}
Diagnostics can be disabled via \texttt{BAXBENCH\_K8S\_DIAGNOSTICS=0}; state
which setting your experiments used.


\subsection{Outcome and Feedback Aggregation}
\label{sec:methods-outcome}

The outcome phase (\texttt{06-outcome/}) closes the iteration.

\paragraph{Feedback artifact.}
\texttt{iteration\_feedback.json} compresses bench and diagnostic signals into
agent-consumable categories: load-test summary, per-endpoint breakdown,
resource pressure (min/avg/max CPU and memory), deployment context, and
failure metadata when applicable.  A human-readable \texttt{.txt} sibling and
append-only \texttt{experiment\_summary.md} support case-study analysis.

\paragraph{Lineage.}
Only \textbf{successful} iterations advance the ``latest good'' code/spec
lineage and become \texttt{prior\_feedback} for the next round.  Failed
iterations receive folder suffixes (\texttt{*-code-failed},
\texttt{*-spec-failed}, \texttt{deploy-failed}) and are excluded from lineage
while remaining on disk for audit.


% ----- §3.5  Failure handling -----
\section{Failure Handling and Iteration Continuity}
\label{sec:methods-failure}

Table~\ref{tab:failure-routing} specifies how failure at iteration $N{-}1$
affects iteration $N$.

\begin{table}[t]
  \centering
  \caption{Iteration routing after failure.}
  \label{tab:failure-routing}
  \small
  \begin{tabular}{@{}lll@{}}
    \toprule
    \textbf{Failed stage} & \textbf{Folder label} & \textbf{Next iteration} \\
    \midrule
    Code (baseline) & \texttt{code-failed} & Sample abort (after max attempts) \\
    Code (refinement) & \texttt{*-code-failed} & Force \texttt{code} \\
    Spec (baseline) & retry in place & Sample abort if exhausted \\
    Spec (refinement) & \texttt{*-spec-failed} & Force \texttt{deployment} \\
    Deploy & \texttt{deploy-failed} & LLM decides \\
  \end{tabular}
\end{table}

\paragraph{Baseline abort.}
If iteration~000 cannot produce passing code or an acceptable spec, refinement
is meaningless; the framework aborts the sample.

\paragraph{Namespace cleanup.}
Before each deploy and after each benchmark, all \texttt{baxbench-*} namespaces
are deleted.  This prevents resource leakage on a fixed lab cluster while
preserving complete on-disk experiment records.


% ----- §3.6  Design principles -----
\section{Design Principles and Constraints}
\label{sec:methods-design}

\paragraph{Bounded search space.}
Raw Kubernetes YAML is vast and error-prone.  Restricting the agent to
\texttt{spec.yaml} keeps comparisons focused on deployment \emph{decisions}
rather than syntactic accidents.

\paragraph{Framework owns infrastructure.}
Network wiring, namespace lifecycle, manifest rendering, Locust orchestration,
and validation rules are fixed.  The agent never executes shell commands on
cluster nodes.

\paragraph{Correctness before performance.}
Functional tests gate all new images; deploy probes gate specs; goodput is
measured only on correct, ready deployments.

\paragraph{Learn from feedback, not from hand-tuned ranges.}
Spec prompts omit recommended numeric values, supporting claims about autonomous
deployment optimisation from measured signals.

\paragraph{Fail-fast refinement, generous baseline.}
Baseline retries bootstrap a working system; single-attempt refinement preserves
a fixed iteration budget and clear per-iteration attribution.

\paragraph{Reproducibility.}
Every LLM call is logged; iteration directories are self-contained experiment
records; cluster topology is codified in \texttt{k8s\_bench/cluster/profiles.py}
rather than ad-hoc host strings in prompts.


% ----- §3.7  Metrics -----
\section{Metrics and Evaluation Criteria}
\label{sec:methods-metrics}

\paragraph{Primary metric.}
\textbf{Sustained goodput} (successful RPS under the adaptive profile), compared
across iterations and scenarios.

\paragraph{Secondary metrics.}
\begin{itemize}
  \item p50 and p95 latency per endpoint;
  \item failure rate during benchmark;
  \item peak step goodput and adaptive-ramp trajectory;
  \item CPU and memory utilisation from \texttt{kubectl top};
  \item deployment parameters in \texttt{spec.yaml} (replicas, requests, placement);
  \item iteration outcome labels and agent decision rationales;
  \item estimated LLM cost per experiment (\texttt{llm\_cost\_ledger.json}).
\end{itemize}

\paragraph{Baselines and ablations (Results chapter).}
Examples for comparison:
\begin{itemize}
  \item baseline deployment only (iteration~000 goodput, no refinement);
  \item forced \texttt{--k8s-refinement deployment} vs.\ \texttt{code} vs.\
        \texttt{auto};
  \item single-shot deploy without iterative feedback.
\end{itemize}


% ----- §3.8  Implementation -----
\section{Implementation Summary}
\label{sec:methods-implementation}

The extension is invoked via \texttt{python src/main.py --mode k8s-bench} or the
wrapper script \texttt{scripts/bench\_k8s.sh}.  Key packages:
\begin{itemize}
  \item \texttt{k8s\_bench/} --- orchestration, stages, cluster profiles,
        spec validation, failure recording;
  \item \texttt{locust\_bench/} --- distributed Locust orchestration and load
        profile registry;
  \item \texttt{bench\_diagnostics/} --- telemetry collectors during benchmark.
\end{itemize}

The canonical iteration pipeline is implemented in
\texttt{k8s\_bench/orchestration/execute.py}.  Design rationale and CLI
reference are documented in \texttt{docs/k8s\_approach.md} and
\texttt{docs/locust\_pipeline.md}.  On-disk layout for a single iteration:

\begin{verbatim}
iteration-NNN-<kind>/
  01-decision/     (refinement only)
  02-code/
  03-spec/
  04-deploy/
  05-bench/
  06-outcome/
  meta.json
  iteration.log
\end{verbatim}

Detailed parameter tables and prompt templates are deferred to the appendix
and repository documentation to keep this chapter focused on experimental
methodology.
